%! Sets of Vectors and the Road to Subspaces — Unit 3, Lesson 3.0 — Homework (Answer Key)
\documentclass[10pt]{article}
\usepackage{linalg-article}
\usepackage{linalg-key}   % pulls in -boxes; do NOT also load -boxes
\newcommand{\TallMath}[1]{$\displaystyle #1\rule[-1.4em]{0pt}{3.2em}$}
\newcommand{\vv}{\mathbf{v}}
\newcommand{\zero}{\mathbf{0}}

\begin{document}

\pageheader{Unit 3, Lesson 3.0}{Homework --- Answer Key}
\namedateperiod

\vspace{0.12in}

\begin{notesbox}{Practice}
For every ``is it a subspace?'' question, run the \textbf{closure test} (adding, scaling, and
$\zero$) and \emph{justify} your answer.

\begin{enumerate}[leftmargin=1.9em, itemsep=12pt]
  \item \textbf{Describe the span} as a \emph{line} or \emph{plane} through the origin, or
        \emph{all of $\mathbb{R}^2$}.
    \begin{enumerate}[label=(\alph*), leftmargin=1.8em, itemsep=4pt]
      \item $\begin{bmatrix} 1 \\ 1 \end{bmatrix}$ and $\begin{bmatrix} 2 \\ 2 \end{bmatrix}$ \hfill \ans{line $y=x$ (parallel: $\begin{bsmallmatrix} 2\\2 \end{bsmallmatrix}=2\begin{bsmallmatrix} 1\\1 \end{bsmallmatrix}$)}
      \item $\begin{bmatrix} 1 \\ 1 \end{bmatrix}$ and $\begin{bmatrix} 1 \\ -1 \end{bmatrix}$ \hfill \ans{all of $\mathbb{R}^2$ (independent)}
      \item $\begin{bmatrix} 3 \\ 6 \end{bmatrix}$ and $\begin{bmatrix} 1 \\ 2 \end{bmatrix}$ \hfill \ans{line $y=2x$ (parallel)}
    \end{enumerate}
  \item \textbf{Subspace or not?} Decide for each set in $\mathbb{R}^2$; if not, give a member and a
        move (add or scale) that escapes it.
    \begin{enumerate}[label=(\alph*), leftmargin=1.8em, itemsep=6pt]
      \item all multiples of $\begin{bmatrix} 2 \\ -1 \end{bmatrix}$ \quad \ans{Subspace --- a line through the origin; closed under $+$ and scaling.}
      \item all points on the line $y = 2x - 1$ \quad \ans{Not --- misses $\zero$ ($x=0\Rightarrow y=-1$). Escape: $2\begin{bsmallmatrix} 0\\-1 \end{bsmallmatrix}=\begin{bsmallmatrix} 0\\-2 \end{bsmallmatrix}$, off the line.}
      \item all $\begin{bmatrix} x \\ y \end{bmatrix}$ with $x = y$ \quad \ans{Subspace --- the line $y=x$ through the origin; sums and scalings keep $x=y$.}
      \item all $\begin{bmatrix} x \\ y \end{bmatrix}$ with $x + y = 1$ \quad \ans{Not --- misses $\zero$ ($0+0\ne 1$). Escape: $2\begin{bsmallmatrix} 1\\0 \end{bsmallmatrix}=\begin{bsmallmatrix} 2\\0 \end{bsmallmatrix}$, and $2+0\ne 1$.}
    \end{enumerate}
  \item \textbf{Column space.} For each matrix, describe $C(A)$ as a set and say whether it is a
        subspace (it always is --- say \emph{why} in one line).
    \begin{enumerate}[label=(\alph*), leftmargin=1.8em, itemsep=4pt]
      \item $A=\begin{bmatrix} 1 & 3 \\ 2 & 6 \end{bmatrix}$ \hfill \ans{line $y=2x$ (columns parallel); a subspace --- a span is always closed.}
      \item $A=\begin{bmatrix} 1 & 0 \\ 0 & 1 \end{bmatrix}$ \hfill \ans{all of $\mathbb{R}^2$ (columns independent); a subspace.}
    \end{enumerate}
  \item \textbf{Explain.} Why must \emph{every} subspace contain the zero vector $\zero$? Answer
        using the scaling rule. \ansline{A subspace has some member $\vv$ and is closed under scaling, so $0\cdot\vv=\zero$ must be in it.}
\end{enumerate}
\end{notesbox}

\begin{extensionbox}
\textbf{A nullspace is a subspace.} Let $A=\begin{bmatrix} 1 & 3 \\ 2 & 6 \end{bmatrix}$. The
\emph{nullspace} is all $\mathbf x$ with $A\mathbf x=\zero$; both rows give $x_1+3x_2=0$.
\begin{enumerate}[label=(\alph*), leftmargin=1.8em, itemsep=5pt]
  \item Find two different nonzero solutions, then describe \emph{all} solutions.
  \item Add your two solutions and check the sum is still a solution --- so the nullspace is
        \emph{closed}, hence a subspace. What line (through the origin?) is it?
\end{enumerate}
\ansline{(a) $x_2=1\Rightarrow\begin{bsmallmatrix} -3\\1 \end{bsmallmatrix}$; $x_2=2\Rightarrow\begin{bsmallmatrix} -6\\2 \end{bsmallmatrix}$. All solutions: $t\begin{bsmallmatrix} -3\\1 \end{bsmallmatrix}$.}
\ansline{(b) Sum $\begin{bsmallmatrix} -3\\1 \end{bsmallmatrix}+\begin{bsmallmatrix} -6\\2 \end{bsmallmatrix}=\begin{bsmallmatrix} -9\\3 \end{bsmallmatrix}$, and $-9+3(3)=0$. \checkmark\ The nullspace is the line through the origin in direction $\begin{bsmallmatrix} -3\\1 \end{bsmallmatrix}$ --- a subspace.}
\end{extensionbox}

\begin{spiralbox}
\textbf{Coming next --- Lesson 3.1: Vector Spaces and Subspaces.} Today you tested sets with the
two closure rules. Next we make those rules precise, list the properties that make $\mathbb{R}^n$
a genuine \emph{vector space}, and prove that spans and column spaces always pass the test.
\end{spiralbox}

\end{document}
